% Propietario conservado: /Users/ruben/Documents/ChatGPT/jueces y controles/REVISION_ARGUMENTAL_APERTURAS_CIERRES_20260915/delivery/ARTICLE_V_ES_ARGUMENTAL_20260915/manuscrito/sections/70_radiacion_termica.tex
% RESULTADO_RECUPERADO: c40_termicas.tex del integral activo.
% Ampliación de prueba: conteo de modos, momentos térmicos, máximos
% espectrales y sustitución de la escala del calendario.
% Las hipótesis de realización se conservan explícitas.
\section{Thermal radiation and the information scale of the calendar}
\label{v:v:sec:radiacion-termica}

Occupation statistics provide the link between the modes of
a realization and their mean energy. In the electromagnetic realization,
the dispersion relation and polarization multiplicity additionally determine the
spectral density. Their composition yields the thermal moments
and spectral maxima, while keeping three operations distinct:
production of the action scale, representation of the modes, and
evaluation of the thermal state.


% BEGIN OWNER /Users/ruben/Documents/ChatGPT/jueces y controles/REVISION_ARGUMENTAL_APERTURAS_CIERRES_20260915/delivery/ARTICLE_V_ES_ARGUMENTAL_20260915/manuscrito/sections/69_enlace_constitutivo_velocidad.tex
% RESULTADO_RECUPERADO: funtor dimensional integral, eq:cZ-internos,
% eq:longitud-ruta y definición de ell_0; composición con artículo III.
% La prueba reunida y su control algebraico son específicos de esta contribución.
\subsection{Constitutive speed and kinematic realization of the clock}
\label{v:v:sec:enlace-velocidad-constitutiva}

Radiative propagation uses the velocity section produced by the
angular transport of the same HMT state. The identification is established
before choosing publication units. Let $x=A_{\mathrm{rad}}$ and
$y=C^*_{\mathrm{rad}}$ be the angular coordinates generated in
section~\ref{v:v:ant:sec:ii-angulos}, with the channels
\eqref{v:v:ant:eq:ii-canales}, and let
$\mathsf R$ be the self-adjoint involution of the two sheets, with rank-one
projectors $P_+$ and $P_-$. On the domain $x>|y|$, the transport and its
channels are
\begin{equation}
 T=\exp(-xI-y\mathsf R)=q_+P_++q_-P_-,
 \qquad q_+=e^{-x-y},\quad q_-=e^{-x+y},\quad 0<q_\pm<1.
 \label{v:v:eq:velocidad-transporte}
\end{equation}
The response of the temporal blocks of lengths $90$ and $120$ is
\begin{equation}
 \mathsf V=(I-T^{90})(I-T^{120})^{-1}
       =r_+P_++r_-P_-,\qquad
 r_\pm=\frac{1-q_\pm^{90}}{1-q_\pm^{120}}.
 \label{v:v:eq:velocidad-respuesta}
\end{equation}
The construction starts from the inherited channels; the symbols $x,y$ and
$q_\pm$ denote their coordinates, not parameters fitted to a target
speed.

\begin{proposition}[Identity of the constitutive and kinematic realizations]
\label{v:v:prop:identidad-velocidad}
In the same dimensional chart, the clock section $c_{\mathrm{int}}$
and the constitutive-response section $c_{\mathrm{vac}}$ coincide:
\begin{equation}
 c_{\mathrm{int}}=c_{\mathrm{vac}}
   =\widehat c\,u_C,
 \qquad
 \widehat c=(r_+r_-)^{-1}
            =\exp(-\mathcal E),
 \label{v:v:eq:identidad-velocidad}
\end{equation}
where
\begin{equation}
 \mathcal E=
 \log\!\bigl[(1-q_+^{90})(1-q_-^{90})\bigr]
 -\log\!\bigl[(1-q_+^{120})(1-q_-^{120})\bigr].
 \label{v:v:eq:velocidad-lector-simetrico}
\end{equation}
The equality preserves sheet exchange.
\end{proposition}
\begin{proof}
The factors appearing in the logarithms are positive. Therefore,
\[
 \mathcal E=\log(r_+r_-)
            =\log\det\mathsf V,
 \qquad \exp(-\mathcal E)=(\det\mathsf V)^{-1}.
\]
The determinant is computed here in the two-sheet realization with rank-one
projectors. The kinematic definition $c_{\mathrm{int}}=\exp(-\mathcal E)u_C$
and the constitutive definition $c_{\mathrm{vac}}=(r_+r_-)^{-1}u_C$
therefore give the same section. Exchanging $P_+\leftrightarrow
P_-$ interchanges $r_+$ and $r_-$ and preserves their product. Finally,
$\det T=e^{-2x}$ corresponds to the elementary transport, whereas
$\det\mathsf V=r_+r_-$ corresponds to its temporal response; the two
determinants have distinct roles in this composition.
\end{proof}

\subsubsection{Path length and elementary duration}

The positive dimensional frame $(u_S,u_C,u_T,u_Q)$ induces the length
basis $u_L=u_Cu_T$. For an enriched path $\gamma$ of
$n=|\gamma|$ transitions and a constant channel, the additive readers are
\begin{equation}
 T_{\mathrm{int}}(\gamma)=n u_T,
 \qquad L_{\mathrm{int}}(\gamma)=n\widehat c\,u_L.
 \label{v:v:eq:velocidad-longitud-reloj}
\end{equation}
Their additivity concerns the count and realization of each transition:
enriched routes additionally preserve orientation and memory.

\begin{corollary}[Normalization of the elementary length]
\label{v:v:cor:longitud-elemental-constitutiva}
For $n>0$, $L_{\mathrm{int}}(\gamma)/T_{\mathrm{int}}(\gamma)
=c_{\mathrm{vac}}$. In the normalization $t_0=u_T$,
\begin{equation}
 \ell_0=c_{\mathrm{int}}t_0=\widehat c\,u_L.
 \label{v:v:eq:longitud-elemental-constitutiva}
\end{equation}
In a general temporal chart $t_0=\tau_0u_T$, with $\tau_0>0$,
the corresponding expression is $\ell_0=\tau_0\widehat c\,u_L$.
\end{corollary}
\begin{proof}
Dividing the expressions in \eqref{v:v:eq:velocidad-longitud-reloj} cancels the integer
$n$ and uses $u_L/u_T=u_C$. The elementary formula is obtained
by multiplying the speed section by $t_0$.
\end{proof}

Generalization to an edge-dependent channel preserves the additive
formula $L_{\mathrm{int}}(\gamma)=\sum_{e\in\gamma}\widehat c(e)u_L$.
Its quotient by $nu_T$ is the mean of the speeds of those edges.
The homogeneous radiative realization uses the constant channel specified
in \eqref{v:v:eq:velocidad-longitud-reloj}.

\subsubsection{Change of chart and adapted normalization}

\begin{proposition}[Covariance of the identification]
\label{v:v:prop:covariancia-identidad-velocidad}
Let $u_C'=d u_C$ and $u_T'=\eta u_T$, with $d,\eta>0$.
The same speed and the same elementary duration have coordinates
\begin{equation}
 \widehat c'=d^{-1}\widehat c,
 \qquad \tau_0'=\eta^{-1}\tau_0,
 \qquad u_L'=d\eta u_L.
 \label{v:v:eq:velocidad-cambio-carta}
\end{equation}
In particular, if $c_V=\widehat c_Vu_C^V$ and
$u_C^V=d u_C^{III}$, equality of the sections is exactly equivalent to
$d\widehat c_V=\widehat c_{III}$.
\end{proposition}
\begin{proof}
Express each section in the two bases. This gives
$\widehat c'u_C'=\widehat c u_C$ and
$\tau_0'u_T'=\tau_0u_T$. Their product satisfies
\[
 \tau_0'\widehat c' u_L'
 =\frac{\tau_0}{\eta}\frac{\widehat c}{d}
       d\eta u_L
 =\tau_0\widehat c u_L=\ell_0.
\]
The final equivalence follows by substituting the relation between the bases.
\end{proof}

The adapted choice $d=\widehat c$ gives $u_C'=c_{\mathrm{vac}}$ and
$\widehat c'=1$. This choice expresses the already constructed section in a
basis adapted to it. The spectral coefficient in the internal chart remains
$(r_+r_-)^{-1}$; spectral inversion uses that internal chart.
For example, retaining the action and charge bases, the constitutive
coordinates transform as
\[
 \widehat\mu'=d\widehat\mu,
 \qquad \widehat\varepsilon'=d\widehat\varepsilon,
 \qquad
 \widehat\mu'\widehat\varepsilon'=(\widehat c')^{-2}.
\]
In the adapted chart, $\widehat\mu'=r_-/r_+$ and
$\widehat\varepsilon'=r_+/r_-$, since the internal coordinates are
$\widehat\mu=r_-^2$ and $\widehat\varepsilon=r_+^2$.

\subsubsection{Transport to the radiative sector}

The dispersion relation of the homogeneous radiative realization is written
$\omega_{\symbf{k}}=c_{\mathrm{vac}}|\symbf{k}|$.
With $h=2\pi\hbar$ and $\beta=(k_B\Theta)^{-1}$, its spectral density
and Stefan--Boltzmann coefficient have the expressions
\begin{equation}
 u_\omega=
 \frac{\hbar\omega^3}{\pi^2(\widehat c u_C)^3}
 \frac{1}{e^{\beta\hbar\omega}-1},
 \qquad
 \sigma=\frac{\pi^2k_B^4}{60\hbar^3(\widehat c u_C)^2}.
 \label{v:v:eq:radiacion-velocidad-constitutiva}
\end{equation}
These substitutions transport precisely the same speed
section. The conditions of the radiative realization are retained:
a free three-dimensional bosonic sector, two transverse polarizations,
zero chemical potential, and positive temperature. Identity
\eqref{v:v:eq:identidad-velocidad} establishes its constitutive antecedent;
mode counting and the thermodynamic limit belong to the specific
proofs for that realization.

\paragraph{Normalization under amplification.}
If refinement acts by $\mathsf V_m=\mathsf V\otimes I_{9^m}$,
the normalized trace $\tau_m=\operatorname{Tr}/(2\cdot9^m)$ satisfies
\[
 \exp[-2\tau_m(\log\mathsf V_m)]
 =\exp[-\log r_+-\log r_-]=\widehat c.
\]
Indeed, each eigenvalue appears $9^m$ times. The multiplicity
cancels in the normalized trace. The ordinary amplified determinant is
$(r_+r_-)^{9^m}$: the two-sheet normalization, or equivalently the
normalized trace, preserves the constitutive reader under that
refinement. This amplification of the representation preserves the
enriched transport from which it arises.

% END OWNER


\subsection{Action, calendar, and temperature}

The clock orientation \(a\) is retained and we abbreviate
\(\hbar=\hbar_a\), \(c=c_{\rm int}=c_{\rm vac}\), with the
identification of \eqref{v:v:eq:identidad-velocidad} in the same chart.
The energy \(\varepsilon_{108}\) of this section is
\(\epsilon_a\) of \eqref{v:v:ant:eq:ii-kb-aut-energia};
\(\Theta_{108}=\Theta_{\rm clk,a}\). Dividing it by \(\log3\)
produces the energy per nat of \eqref{v:v:ant:eq:ii-kb-aut-energia-nat}.
The construction and covariance of the thermal pair are given in
section~\ref{v:v:ant:sec:ii-kb-aut-desarrollo}.

We write \(h=2\pi\hbar>0\) for the action publication, \(t_0>0\)
for the time unit of its dimensional chart and \(\ell_0=ct_0\) for
the corresponding length. The chronological cell of length \(108\)
determines the energy
\begin{equation}
 \varepsilon_{108}=\frac{h}{108t_0},\qquad
 k_B\Theta_{108}\log 3=\varepsilon_{108},\qquad
 \beta_{108}=\frac{108t_0\log3}{h}.
 \label{v:v:eq:energia-informacion-calendario}
\end{equation}
Here \(\beta=(k_B\Theta)^{-1}\) is the Gibbs parameter. The equality
fixes the energy-valued product of temperature and the entropy
transducer. The dimensional specification of \(k_B\) and
\(\Theta_{108}\) separately uses the thermal section and its bases. The natural
logarithm measures information per tritic branch; its factor \(\log3\) converts
the cardinality of the cell into dimensionless entropy.

Also introduce the associated units
\begin{equation}
 t_P=\frac{108}{2\pi}t_0,\qquad
 \ell_P=ct_P,\qquad E_P=\frac{\hbar}{t_P}
 =\varepsilon_{108}.
 \label{v:v:eq:unidades-termicas-calendario}
\end{equation}
These equalities specify the meaning of the symbols in this section;
their evaluation in another chart preserves the identities between quotients.
The length \(\ell_P=(54/\pi)\ell_0\) coincides with the areal unit
of section~\ref{v:v:ant:sec:gravity}, in the circular realization
declared there.

\subsection{Mode density and spectral law}

\begin{proposition}[Thermal realization of transverse modes]
\label{v:v:prop:ley-espectral}
Consider a free realization in a three-dimensional periodic box,
whose excitation space contains the wave vectors
\(\mathbf k\ne0\), with two transverse modes per vector, dispersion
\(\omega=c|\mathbf k|\), excitation energy \(\hbar\omega\),
zero chemical potential, and \(\beta>0\). The constant mode is kept
fixed as background data. In the thermodynamic limit, the energy density
per angular frequency is
\begin{equation}
 u_\omega(\Theta)=\frac{\hbar\omega^3}{\pi^2c^3}
 \frac{1}{e^{\beta\hbar\omega}-1},\qquad \omega>0.
 \label{v:v:eq:planck-angular}
\end{equation}
The energy is that of the excitations above the vacuum of this
representation.
\end{proposition}

\begin{proof}
For a box of volume \(V\), the wave vectors have density
\(V/(2\pi)^3\). In a radial shell of thickness \(dk\), the density
per unit volume, including the two polarizations, is
\[
 2\frac{4\pi k^2\,dk}{(2\pi)^3}=\frac{k^2}{\pi^2}\,dk
 =\frac{\omega^2}{\pi^2c^3}\,d\omega.
\]
This is the limiting density of the sums over the periodic lattice
with the constant mode omitted. This choice defines the partition product
before taking the limit: each excitation mode has positive energy.
Excluding a single point preserves the limiting radial density.
The previously obtained bosonic occupation is
\(\bar n_\omega=(e^{\beta\hbar\omega}-1)^{-1}\).
Multiplying the density by \(\hbar\omega\bar n_\omega\) proves
\eqref{v:v:eq:planck-angular}. Near zero, the energy integrand is
of order \(\omega^2\), and at infinity it decays exponentially, so
its integral converges.
\end{proof}

The density arises from the specified three-dimensional representation.
The transverse multiplicity used here has that physical domain;
the dimension of a harmonic space on the APP torus describes,
separately, the cohomology of its cellular complex.


% BEGIN OWNER /Users/ruben/Documents/ChatGPT/jueces y controles/REVISION_ARGUMENTAL_APERTURAS_CIERRES_20260915/delivery/ARTICLE_V_ES_ARGUMENTAL_20260915/manuscrito/sections/71_limite_termodinamico.tex
% FORMALIZACION_NUEVA de la justificación del límite utilizado en
% 70_radiacion_termica.tex. No introduce un generador de constantes.
\subsection{Thermodynamic limit of occupation sums}
\label{v:v:subsec:limite-termodinamico}

The preceding radial density can be obtained as an effective limit of the
sums of the periodic realization. The estimate must preserve the fixed
constant mode and uniformly control both the origin and the spectral
tail. The action and calendar publications remain prior
to this realization operation.

\begin{lemma}[Convergence of the three thermal observables]
\label{v:v:lem:limite-termodinamico}
Let \(a,b>0\), \(L>0\), \(\Delta=2\pi/L\), and
\[
 F_N(r)=\frac{1}{e^{ar}-1},\qquad
 F_E(r)=\frac{br}{e^{ar}-1},\qquad
 F_Z(r)=-\log(1-e^{-ar})\quad(r>0).
\]
For each of these functions,
\begin{equation}
 \lim_{L\to\infty}\frac{2}{L^3}
 \sum_{\nu\in\mathbb Z^3\setminus\{0\}}
 F\!\left(\frac{2\pi}{L}|\nu|\right)
 =\frac{1}{\pi^2}\int_0^\infty r^2F(r)\,dr.
 \label{v:v:eq:limite-modos-discretos}
\end{equation}
The sums and the integral are finite. The specialization
\(a=\beta\hbar c\), \(b=\hbar c\) gives, respectively,
the number density, excitation-energy density, and logarithmic partition-function
density of the transverse modes.
\end{lemma}

\begin{proof}
All three functions are continuous and positive on \((0,\infty)\).
There exist constants \(C_0,C_1,d>0\), depending on \(a,b\) but
independent of \(\Delta\), such that
\begin{equation}
 F(r)\leq\frac{C_0}{r}\quad(0<r\leq1),\qquad
 F(r)\leq C_1e^{-dr}\quad(r\geq1).
 \label{v:v:eq:envolventes-termicas}
\end{equation}
Indeed, \(e^{ar}-1\geq ar\) bounds \(F_N\) and \(F_E\)
near the origin; moreover,
\(F_Z(r)=\log(1+F_N(r))\leq F_N(r)\leq1/(ar)\).
For \(r\geq1\), the inequality
\(-\log(1-y)\leq y/(1-y)\), with \(y=e^{-ar}\), gives the
exponential bound for \(F_Z\). The other two follow from
\((e^{ar}-1)^{-1}\leq e^{-ar}/(1-e^{-a})\), absorbing the
factor \(r\) in \(F_E\) into an exponential of smaller decay rate. These bounds
already prove integrability of \(r^2F(r)\).

For the discrete sums, group the lattice points by
\(m=\|\nu\|_\infty\geq1\). Each shell contains exactly
\[
 (2m+1)^3-(2m-1)^3=24m^2+2
\]
points and satisfies \(m\leq|\nu|\leq\sqrt3m\).
Fix \(0<\delta\leq1\) and \(0<\Delta\leq1\).
If \(\delta<\Delta\), there are no points with
\(0<\Delta|\nu|\leq\delta\). Otherwise, setting
\(N=\lfloor\delta/\Delta\rfloor\), the first bound in
\eqref{v:v:eq:envolventes-termicas} gives
\begin{align*}
 \Delta^3\!\sum_{0<\Delta|\nu|\leq\delta}F(\Delta|\nu|)
 &\leq C_0\Delta^2\sum_{m=1}^{N}\frac{24m^2+2}{m}\\
 &\leq26C_0\Delta^2\sum_{m=1}^{N}m
 \leq26C_0\delta^2.
\end{align*}
The second line uses \(m^{-1}\leq m\) and then
\(N(N+1)/2\leq N^2\) for \(N\geq1\).
The integral contribution from the same neighborhood is also of order
\(\delta^2\), since
\(\int_0^\delta r^2F(r)\,dr\leq C_0\delta^2/2\).

For \(R\geq1\), every point with \(\Delta|\nu|>R\) belongs
to a shell \(m>R/(\sqrt3\Delta)\). The second bound in
\eqref{v:v:eq:envolventes-termicas} gives
\begin{align*}
 \Delta^3\!\sum_{\Delta|\nu|>R}F(\Delta|\nu|)
 &\leq C_1e^{-dR/(2\sqrt3)}
 \Delta^3\sum_{m\geq1}(24m^2+2)e^{-d\Delta m/2}\\
 &\leq C_2e^{-dR/(2\sqrt3)},
\end{align*}
where \(C_2\) is independent of \(0<\Delta\leq1\). To verify
this uniformity, set \(q=e^{-d\Delta/2}\) and use
\[
 \sum_{m\geq1}q^m=\frac{q}{1-q},\qquad
 \sum_{m\geq1}m^2q^m=\frac{q(1+q)}{(1-q)^3}.
\]
The first identity is the geometric series, and the second follows by
applying \(q\,d/dq\) twice, which is valid for \(|q|<1\).
Moreover,
\(1-e^{-d\Delta/2}\geq\Delta(1-e^{-d/2})\) by concavity.
The factor \(\Delta^3\) therefore compensates for the denominators of
both expressions. The integral tail likewise tends to zero by the
exponential bound. For each \(\Delta>0\), the same estimates
prove convergence of the infinite sum.

On the compact annulus \(\delta\leq|k|\leq R\), the function
\(F(|k|)\) is continuous and bounded. Riemann sums of mesh
\(\Delta\) converge to its integral; the two boundary spheres
have measure zero. First let \(\Delta\to0\), then
\(\delta\to0\), and finally \(R\to\infty\), using the preceding
uniform bounds. The result is
\[
 \Delta^3\sum_{\nu\ne0}F(\Delta|\nu|)
 \longrightarrow \int_{\mathbb R^3}F(|k|)\,d^3k
 =4\pi\int_0^\infty r^2F(r)\,dr.
\]
Since \(2/L^3=2\Delta^3/(2\pi)^3\), the final radial coefficient
is \(2\cdot4\pi/(2\pi)^3=1/\pi^2\), proving
\eqref{v:v:eq:limite-modos-discretos}.
\end{proof}

Control at the origin is carried out over positive-energy excitations.
The constant mode remains fixed before the partition function is formed:
removing it afterward from an integral would not replace that condition on
the finite representation. The proof establishes the limit of the
specified realization and preserves the distinction between that realization and
the prior construction of its action and calendar parameters.

% END OWNER


\subsection{Thermal moments and information}

\begin{lemma}[Bosonic moment]
\label{v:v:lem:momento-bosonico}
For every real \(s>1\),
\begin{equation}
 \int_0^\infty\frac{x^{s-1}}{e^x-1}\,dx
 =\Gamma(s)\zeta(s).
 \label{v:v:eq:momento-bosonico}
\end{equation}
\end{lemma}
\begin{proof}
For \(x>0\), \((e^x-1)^{-1}=\sum_{m\geq1}e^{-mx}\).
All summands are nonnegative; Tonelli's theorem allows the sum
and integral to be interchanged. The substitution \(y=mx\) gives
\(\int_0^\infty x^{s-1}e^{-mx}\,dx=m^{-s}\Gamma(s)\).
The sum converges for \(s>1\), yielding the stated equality.
\end{proof}

\begin{proposition}[Number, energy, and entropy densities]
\label{v:v:prop:densidades-termicas}
In the realization of Proposition~\ref{v:v:prop:ley-espectral},
\begin{align}
 n_\gamma&=\frac{2\zeta(3)}{\pi^2c^3(\beta\hbar)^3},\\
 u_\gamma&=\frac{\pi^2}{15c^3\hbar^3\beta^4},\\
 s_\gamma&=\frac{4\pi^2 k_B}{45c^3(\beta\hbar)^3}.
 \label{v:v:eq:densidades-termicas}
\end{align}
At the calendar temperature, these relations become
\begin{equation}
 n_\gamma\ell_P^3=\frac{2\zeta(3)}{\pi^2(\log3)^3},\quad
 \frac{u_\gamma\ell_P^3}{E_P}=\frac{\pi^2}{15(\log3)^4},\quad
 \frac{s_\gamma\ell_P^3}{k_B}=\frac{4\pi^2}{45(\log3)^3}.
 \label{v:v:eq:densidades-calendario}
\end{equation}
\end{proposition}

\begin{proof}
Integrating occupation against the mode density uses
\eqref{v:v:eq:momento-bosonico} with \(s=3\), and the energy integral
uses it with \(s=4\). This gives \(\Gamma(3)=2\) and
\(\Gamma(4)\zeta(4)=6\pi^4/90=\pi^4/15\).
For entropy, start from the partition product of the modes:
\[
 \frac{\log Z}{V}=-\frac{1}{\pi^2c^3}
 \int_0^\infty\omega^2
 \log(1-e^{-\beta\hbar\omega})\,d\omega
 =\frac{\pi^2}{45c^3(\beta\hbar)^3}.
\]
The last equality follows by integration by parts after the substitution
\(x=\beta\hbar\omega\). The boundary terms vanish:
\(x^3\log(1-e^{-x})\) tends to zero at both endpoints.
The Gibbs identity \(s_\gamma/k_B=\log Z/V+\beta u_\gamma\)
proves the third expression. Finally,
\(\beta_{108}\hbar=t_P\log3\) and \(\ell_P=ct_P\) give the three
calendar normalizations.
\end{proof}

\subsection{Thermal flux and the Stefan--Boltzmann constant}

\begin{proposition}[Isotropic flux and action normalization]
\label{v:v:prop:stefan}
The hemispherical flux of the isotropic realization is
\begin{equation}
 j_*(\Theta)=\sigma\Theta^4,\qquad
 \sigma=\frac{\pi^2 k_B^4}{60\hbar^3c^2}
 =\frac{2\pi^5k_B^4}{15h^3c^2}.
 \label{v:v:eq:stefan}
\end{equation}
At the calendar temperature,
\begin{equation}
 j_*(\Theta_{108})
 =\frac{2\pi^5}{15\,108^4(\log3)^4}
 \frac{h}{\ell_0^2t_0^2}.
 \label{v:v:eq:stefan-calendario}
\end{equation}
\end{proposition}
\begin{proof}
The isotropic energy per unit solid angle is
\(u_\gamma/(4\pi)\). Flux normal to a surface integrates
\(c\cos\theta\) over the hemisphere:
\[
 \frac{cu_\gamma}{4\pi}
 \int_0^{2\pi}\!d\phi\int_0^{\pi/2}\cos\theta\sin\theta\,d\theta
 =\frac{cu_\gamma}{4}.
\]
The formula for \(u_\gamma\), \(\beta^{-1}=k_B\Theta\), and
\(h=2\pi\hbar\) prove \eqref{v:v:eq:stefan}.
Substitution of \eqref{v:v:eq:energia-informacion-calendario} and
\(c=\ell_0/t_0\) gives \eqref{v:v:eq:stefan-calendario}.
\end{proof}

Within this realization, the constant \(\sigma\) is the coefficient
combining bosonic occupation, transverse-mode density, and
angular projection of the flux. The fourth-power dependence arises from the cubic
frequency moment together with the integration differential. The calendar
readout expresses that same flux through action, length,
time, and tritic information.

\newpage
\subsection{Spectral maxima and displacement law}

\begin{lemma}[Positive root for spectral maxima]
\label{v:v:lem:raiz-wien}
For every integer \(n>1\), the equation
\(n(1-e^{-x})=x\) has a unique solution \(x_n>0\).
The function \(f_n(x)=x^n/(e^x-1)\) has its unique positive maximum at
\(x_n\).
\end{lemma}
\begin{proof}
Let \(g_n(x)=n(1-e^{-x})-x\). We have
\(g_n(0)=0\), \(g_n'(x)=ne^{-x}-1\), and
\(g_n''(x)=-ne^{-x}<0\). The maximum of \(g_n\) is attained at
\(\log n\), where its value is \(n-1-\log n>0\), whereas
\(g_n(n)=-ne^{-n}<0\). Strict monotonicity on either side of the maximum
proves existence and uniqueness of the positive root. The sign of
\(f_n'(x)\) agrees with that of \(g_n(x)\); moreover, \(f_n\) tends
to zero at zero and infinity. This proves uniqueness of the maximum.
\end{proof}

\begin{proposition}[Frequency and wavelength readouts]
\label{v:v:prop:wien}
The maxima of radiance per frequency and per wavelength are
\begin{equation}
 \nu_{\max}=\frac{x_3 k_B\Theta}{h},\qquad
 \lambda_{\max}=\frac{hc}{x_5 k_B\Theta}.
 \label{v:v:eq:wien}
\end{equation}
In particular,
\begin{equation}
 \nu_{\max}(\Theta_{108})=
 \frac{x_3}{108t_0\log3},\qquad
 \lambda_{\max}(\Theta_{108})=
 \frac{108\log3}{x_5}\ell_0.
 \label{v:v:eq:wien-calendario}
\end{equation}
\end{proposition}
\begin{proof}
The conversion \(\omega=2\pi\nu\) and the isotropic angular factor give
\[
 B_\nu=\frac{2h\nu^3}{c^2}\frac1{e^{\beta h\nu}-1}.
\]
With \(\nu=c/\lambda\), the transformed density satisfies
\(B_\lambda=B_\nu|d\nu/d\lambda|\); hence
\[
 B_\lambda=\frac{2hc^2}{\lambda^5}
 \frac1{e^{\beta hc/\lambda}-1}.
\]
The respective changes of variable lead to \(f_3\) and \(f_5\).
The preceding lemma proves \eqref{v:v:eq:wien}; the calendar scale
gives \eqref{v:v:eq:wien-calendario}. The density Jacobian explains why
the two maxima correspond to distinct roots:
\(\nu_{\max}\lambda_{\max}=c x_3/x_5\).
\end{proof}

Wien displacement is structurally defined by the stationary point of
the chosen spectral density. A change of variable transforms
both the argument and the measure. This precision locates
each spectral constant in its generating operation and preserves its
physical meaning when passing between frequency and wavelength.
